
\section{Function shift sum}

\bigskip

\begin{lemma}
\label{fun_shift_sum_one}
\lean{fun_shift_sum_one}
\leanok
\uses{def:SF, def:d}
If $f:\Zz^n \to \Fin n$ is \eqref{SF}, $b\in\Zz^n$, and $k\in\Fin n$, then
$$\sum_{i\in\Fin n} \delta(f(b+{\rm e}_i),k) = 1.$$
\end{lemma}

\begin{proof}
\uses{fun_shift_inj}
Since $f$ is \eqref{SF}, we {\B obtain} $i_k\in\Fin n$ such that $f(b+{\rm e}_{i_k}) = k$.

We shall {\B rewrite} the sum by proving that {\B only the $i_k$ term in the sum is equal to $1$, and the rest are $0$}.

Clearly, $\delta(f(b+{\rm e}_{i_k}),k) = 1$, in view of \eqref{d}.

{\B Introduce} $i\in\Fin n$ with $i\neq i_k$. {\B Rewritting} $k$ as  $f(b+{\rm e}_{i_k})$ and {\B applying} the logic of {\B if-then-else} in \eqref{d}, we have to prove that $f(b+{\rm e}_i) \neq f(b+{\rm e}_{i_k})$. In order to get a {\B contradiction}, assume that $f(b+{\rm e}_i) = f(b+{\rm e}_{i_k})$. Then, by {\B Lemma} \ref{fun_shift_inj}, we {\B have} $i=i_k$, which is {\B false} in view of the assumption $i\neq i_k$.
\end{proof}

\verb|Mathlib hints:| \href{https://leanprover-community.github.io/mathlib4_docs/Mathlib/Data/Fintype/BigOperators.html#Fintype.sum_eq_single}{rw [Fintype.sum\_eq\_single $i_k$]}, \href{https://leanprover-community.github.io/mathlib4_docs/Init/Core.html#if_pos}{if\_pos}, \href{https://leanprover-community.github.io/mathlib4_docs/Init/Core.html#if_neg}{if\_neg}.

\bigskip

\begin{lemma}
\label{fun_shift_sum_pow2}
\lean{fun_shift_sum_pow2}
\leanok
\uses{def:SF, def:d}
 If $f:\Zz^n \to \Fin n$ is \eqref{SF} and $k\in\Fin n$, then
$$\sum_{b\in\Zz^n} \sum_{i\in\Fin n} \delta(f(b+{\rm e}_i),k) = 2^n.$$
\end{lemma}

\begin{proof}
\uses{fun_shift_sum_one}
By  {\B Lemma \ref{fun_shift_sum_one}}, we {\B have}:
$$\forall b\in\Zz^n, \sum_{i\in\Fin n} \delta(f(b+{\rm e}_i),k) = 1.$$ With {\B only} this relation, we can {\B simplify}
$\displaystyle\sum_{b\in\Zz^n} \sum_{i\in\Fin n} \delta(f(b+{\rm e}_i),k)$ to  $\displaystyle\sum_{b\in\Zz^n} 1$,
which {\B simplifies} further to $2^n$.
\end{proof}

\bigskip

\begin{lemma}
\label{fun_shift_sum_rev}
\lean{fun_shift_sum_rev}
\leanok
\uses{def:e, def:d}
 If $f:\Zz^n \to \Fin n$ and $k,i\in\Fin n$, then
$$\sum_{b\in\Zz^n} \delta(f(b+{\rm e}_i),k) = \sum_{b\in\Zz^n} \delta(f(b),k).$$
\end{lemma}

\begin{proof}
{\B Let} $ g: \Zz^n \to \Zz^n, g(b) = b + {\rm e}_i$. Then we {\B have} that $g$ is {\B bijective}. {\B Applying} the fact that $g$ is {\B bijective} to our  {\B finite} sum, the equality holds, because $\displaystyle\sum_{b\in\Zz^n} \delta(f(g(b)),k)$ is {\B simply} $\displaystyle\sum_{b\in\Zz^n}\delta(f(b + {\rm e}_i),k).$
\end{proof}

\noindent\verb|Mathlib hints:| \href{https://leanprover-community.github.io/mathlib4_docs/Mathlib/Logic/Function/Defs.html#Function.Bijective}{Function.Bijective}, \href{https://leanprover-community.github.io/mathlib4_docs/Mathlib/Algebra/Group/Defs.html#add_left_injective}{add\_left\_injective}, \href{https://leanprover-community.github.io/mathlib4_docs/Mathlib/Algebra/Group/Basic.html#add_right_surjective}{add\_right\_surjective}, \href{https://leanprover-community.github.io/mathlib4_docs/Mathlib/Algebra/BigOperators/Group/Finset/Defs.html#Fintype.sum_bijective}{Fintype.sum\_bijective}.
